\documentclass[a4paper,12pt]{article}
\input{LaTeX-ToolKit/template.tex}
\begin{document}
\title{Tree} 
\author{沈威宇}
\date{\temtoday}
\titledoc
\sct{Tree}
\ssc{Tree}
\sssc{Tree}
A tree is an abstract data type of a set of nodes. A relationship between two nodes $i,j$ is defined with
\bit
\item $i$ is called a child of $j$ and
\item $j$ is called a parent of $i$; and
\item a node can have at most one parent.
\eit
If not empty, a tree is recursively defined as a distinguished node $r$ called root and zero or more nonempty trees called subtrees, where the root of each substree is a child of $r$.
\sssc{Terminology}
\bit
\item Siblings: Nodes sharing the same parent are called siblings. Children of a node are defined with a left to right order.
\item Leaf: A node with no children is called a leaf, external node, or exterior node.
\item Internal node: A node that is not a leaf is called an internal node or interior node.
\item Depth: The number of edges from the root to a node is called the depth of the node, while the root has depth $0$.
\item Height: The height of a leaf is defined as 0. The height of an internal node is defined as the maximum height of its children plus 1. The maximum depth of the nodes, i.e., the height of its root, is called the height or depth of the tree and denoted $h$. The height or depth of an empty tree is defined as -1. Some define heights of nodes and trees as the definition here plus 1, which is not used here.
\item Ancestors: The ancestors of a node is recursively defined as itself for root and itself and the ancestors of its parent for non-root nodes.
\item Descendants: The descendants of a node is recursively defined as itself for leaves and itself and the descendants of its children for internal nodes.
\item Level: Nodes of the same depth are called a level. Levels are ordered by depth. The left-to-right order of nodes in a level is defined recursively to, if they are siblings, follow that of siblings; otherwise follow that of their parents.
\item Subtree: A tree consisting of all descendants of a node is called a subtree. The left-to-right order of subtrees whose roots are of the same depth are defined as that of their roots. A subtree whose root is a child of a node is called a subtree of the node.
\item The umber of nodes in a tree is denoted $n$.
\eit
\sssc{Tree representation}
A representation of tree where each pair of parent and child is connected with a edge.
\sssc{Left child, right sibling representation}
A representation of tree where each internal node and its leftmost child is connected with a edge, and each sibling that is not the leftmost one is called to the sibling at its left.
\sssc{Skewed}
A tree whose every level contains only one node is called skewed.
\ssc{$n$-ary tree}
\sssc{$n$-ary tree}
A tree in which no node can have more than $n$ children is called an $n$-ary tree. A child is indexed into one and only one integer in $[1,n]$ even if its parent has less than $n$ children.
\sssc{Full, proper, or plane}
An $n$-ary tree where each internal node has exactly $n$ children is called full, proper, or plane.
\sssc{Filled}
A level of depth $k\in\bbN$ of an $n$-ary tree is called filled iff each node in the level of depth $k-1$ has exactly $n$ children. The level of depth 0 is called filled iff the tree is nonempty.
\sssc{Complete}
An $n$-ary tree where every level, except possibly the last, is filled and whose nodes in the last level are as far left as possible, is called complete.

A complete $n$-ary tree is indexed from top to bottom and then from left to right, with the root as node $1$.
\sssc{Perfect}
A complete $n$-ary tree where all leaves are of the same depth is called perfect.

A perfect $n$-ary tree with height $k$ has $\sum_{j=0}^kn^j$ nodes.
\ssc{Binary tree}
\sssc{Binary tree}
In a binary tree, the two possible children of a node are called left child and right child. A child is either left or right but not both even if it has no sibling. The left and right subtrees of a node $x$ are denoted $\operatorname{left}(x)$ and $\operatorname{right}(x)$ respectively.
\sssc{Complete binary tree}
For a complete binary tree with $n$ nodes:
\bit
\item Parent of node $i$ is node $\floor{\frac{i}{2}}$ for all integers $i>1$.
\item The left child of node $i$ is $2i+1$ for all integers $i\le\frac{n-1}{2}$.
\item The right child of node $i$ is $2i$ for all integers $i\le\frac{n}{2}$.
\eit
\sssc{Perfect binary tree}
A perfect binary tree with height $k$ has $2^{k+1}-1$ nodes.
\sssc{Balanced tree}
A type of binary tree with height $O(\log n)$ where $n$ is the number of nodes is called balanced.
\sssc{Balance factor}
The height of the left subtree minus the height of the right subtree of a node in a binary tree is called the balance factor of the node.
\sssc{Left-skewed and right-skewed}
A binary tree whose every internal node has left child and does not have right child is called left-skewed.

A binary tree whose every internal node has right child and does not have left child is called right-skewed.
\ssc{Tree traversal}
They are defined recursively. Suppose the root is $r$.
\sssc{Preorder traversal}
First visit $r$, and then visit the subtrees of $r$ from left to right.
\sssc{Postorder traversal}
Visit the subtrees of $r$ from left to right, and finally visit $r$.
\sssc{Inorder traversal (binary tree)}
First visit the left subtree of $r$, and then visit $r$, and then visit the right subtree of $r$.
\sssc{Full binary tree with distinct values can be uniquely determined by its preorder and postorder traversals}
A full binary tree with distinct values can be uniquely determined by its preorder and postorder traversals.
\bpr
The following C code provides an algorithm to reconstruct the tree recursively:
\lstinputlisting[language=C]{reconstructFullBinaryTreePreorderPostorder.c}
\epr
\sssc{Binary tree with distinct values cannot be uniquely determined by its preorder and postorder traversals}
The statement that a binary tree with distinct values can be uniquely determined by its preorder and postorder traversals is false.
\bpr
No. A counterexample is the following two binary trees
\begin{center}
\begin{forest}
[1[2][,phantom]]
\end{forest}
\begin{forest}
[1[,phantom][2]]
\end{forest}
\end{center}
Both trees have preorder traversal 1, 2, and postorder traversal 2, 1, but they are different.
\epr
\sssc{Binary tree with distinct values can be uniquely determined by its preorder and inorder traversals}
A binary tree with distinct values can be uniquely determined by its preorder and inorder traversals.
\bpr
The following C code provides an algorithm to reconstruct the tree recursively:
\lstinputlisting[language=C]{reconstructBinaryTreePreorderInorder.c}
\epr
\sssc{Binary tree with distinct values can be uniquely determined by its postorder and inorder traversals}
A binary tree with distinct values can be uniquely determined by its postorder and inorder traversals.
\bpr
The following C code provides an algorithm to reconstruct the tree recursively:
\lstinputlisting[language=C]{reconstructBinaryTreePostorderInorder.c}
\epr
\ssc{Binary search tree (BST)}
\sssc{Binary search tree (BST)}
A binary search tree stores keys in the nodes of a binary tree so that searching, insertion and deletion can be done efficiently. For every node $x$, all the keys in its left subtree are smaller than $x$ and all keys in its right subtree are larger than $x$. Thus, the values are distinct.
\sssc{Searching}
Suppose we are searching for $x$ and the root is $r$.
\bit
\item If $x=r$, then we are done.
\item If $x<r$, then we search in the left subtree.
\item If $x>r$, then we search in the right subtree.
\eit
The time complexity is $O(h)$.
\sssc{Finding minimum and maximum}
To find the node containing the smallest key, start at root and go to the left child as long as there is one, and the leaf we stop at is the smallest key. Similarly for find the largest key.

The time complexity is $O(h)$.
\sssc{Insertion}
Suppose we want to insert $x$. First search for it, if we found it, do noting (or update something); otherwise, if we stop at $y$,
\bit
\item if $x>y$, insert $x$ as the right child of $y$ if $x>y$, and
\item if $x<y$, insert $x$ as the left child of $y$.
\eit
The time complexity is $O(h)$.
\sssc{Deletion}
Suppose we want to delete $x$. First search for it, if we don't found it, do noting (or update something); otherwise,
\bit
\item if it is a leaf, delete it,
\item if it has one child $y$, adjust $y$ as a child of the parent of $x$ of the same position as $x$ and delete $x$, and
\item if it has two children, choose either of the policies:
\bit
\item inorder successor: replace the value of $x$ with the value of the minimum element $y$ in the right subtree of $x$, i.e. inorder successor, i.e., the leftmost descendant of the right child of $x$, which has no left child but may have right child, and
\bit
\item remove node $y$ if it has no child, and
\item replace node $y$ with its right child if it has one;
\eit
\item inorder predecessor: replace the value of $x$ with the value of the maximum element $y$ in the left subtree of $x$, i.e. inorder predecessor, i.e., the rightmost descendant of the left child of $x$, which has no right child but may have left child, and
\bit
\item remove node $y$ if it has no child, and
\item replace node $y$ with its left child if it has one;
\eit
\eit
\eit
The time complexity is $O(h)$.
\sssc{(Simple) rotation}
A (simple) rotation on an edge between a node $x$ and its parent $p$, also called around $p$, is a tree operation on BST and has the following two symmetric variants:
\bit
\item Right rotation: This is done when $x$ is a left child of $p$. We perform the following assignments:
\[\operatorname{left}(p)=\operatorname{right}(x)\]
\[\operatorname{right}(x)=p\]
\item Left rotation: This is done when $x$ is a right child of $p$. We perform the following assignments:
\[\operatorname{right}(p)=\operatorname{left}(x)\]
\[\operatorname{left}(x)=p\]
\eit
\sssc{Uniqueness given preorder traversal}
A BST is uniquely determined by its preorder traversal.
\bpr
A BST with one node is obviously unique. The first element must be root, the second element to the last element smaller than root must be the left subtree, and the rest must be the right subtree. Recursively partition into the subtrees. We uniquely obtained the BST.
\epr
\sssc{Uniqueness given postorder traversal}
A BST is uniquely determined by its postorder traversal.
\bpr
A BST with one node is obviously unique. The last element must be root, the first element to the last element smaller than root must be the left subtree, and the rest must be the right subtree. Recursively partition into the subtrees. We uniquely obtained the BST.
\epr
\sssc{Binary tree is BST iff inorder traversal is in ascending order}
A binary tree with distinct values is a binary search tree iff its inorder traversal is sorted in ascending order.
\bpr
Only if: Prove by mathematical induction. The inorder traversal of a BST with one node is sorted in ascending order. Assume the inorder traversal of any BST with fewer than $n$ nodes of distinct values is sorted in ascending order. Consider any BST with $n$ nodes, the inorder traversal visits the left subtree, which is smaller than root and itself sorted in ascending order since it's a BST with $n-1$ nodes, and then root, and then the right subtree, which is larger than root and itself sorted in ascending order since it's a BST with $n-1$ nodes. Thus, we conclude that the inorder traversal of any BST with $n$ nodes is sorted in ascending order. By mathematical induction, the statement holds.

If: Prove by mathematical induction. A binary tree with one node is a BST. Assume that any binary tree with fewer than $n$ nodes of distinct values whose inorder traversal is sorted in ascending order is a BST. Consider any binary tree with $n$ nodes of distinct values and root $r$ whose inorder traversal is sorted in ascending order. The subsequence before $r$ is a binary tree with fewer than $n$ nodes of distinct values whose inorder traversal is sorted in ascending order whose values of nodes are smaller than $r$. The subsequence after $r$ is a binary tree with fewer than $n$ nodes of distinct values whose inorder traversal is sorted in ascending order whose values of nodes are larger than $r$. Since the left subtree is a BST and smaller than root and the right subtree is a BST and larger than root, we conclude that any binary tree with $n$ nodes of distinct values whose inorder traversal is sorted in ascending order is a BST. By mathematical induction, the statement holds.
\epr
\sssc{Number of distinct BST shapes}
The number $T_n$ of distinct shapes of BSTs with $n$ nodes is:
\[T_0=1\]
\[T_n=\sum_{i=0}^{n-1}T_iT_{n-1-i},\quad n\in\bbN\]
\[T_n=\frac{1}{(n+1)}\binom{2n}{n},\quad n\in\bbN_0\]
\bpr
Obviously
\[T_0=T_1=1\]
Let the values be
\[k_1<k_2<\ldots<k_n\]
Suppose we choose $k_i$ as root, where $1\le i\le n$. The left subtree contains $k_1,\ldots,k_{i-1}$, so it has $i-1$ nodes and $T_{i-1}$ possible combinations. The right subtree contains $k_{i+1},\ldots,k_n$, so it has $n-1$ nodes and $T_{n-i}$ possible combinations.

Summing over all choices gives
\[T_n=\sum_{i=1}^nT_{i-1}T_{n-i}=\sum_{i=0}^{n-1}T_iT_{n-1-i}\]
Using generating function:
\[C(x)=\sum_{n=0}^{\infty}T_nx^n\]
\[C(x)=1+xC(x)^2\]
\[C(x)=\frac{1\pm\sqrt{1-4x}}{2x}\]
\[C(0)=\mp\frac{1}{\sqrt{1-4x}}=T(0)=1\]
\[C(x)=\frac{1-\sqrt{1-4x}}{2x}\]
\[\ba
\sqrt{1-4x}&=\sum_{k=0}^{\infty}\binom{0.5}{k}(-4x)^k\\
           &=1+\sum_{k=1}^{\infty}\frac{\prod_{i=0}^{k-1}(0.5-i)}{k!}(-4x)^k\\
           &=1+\sum_{k=1}^{\infty}\frac{2^k\prod_{i=0}^{k-1}(2i-1)}{k!}x^k\\
           &=1+\sum_{k=1}^{\infty}-\frac{2^k\prod_{i=1}^{k-1}(2i-1)}{k!}x^k\\
           &=1+\sum_{k=1}^{\infty}-\frac{2^k\prod_{i=1}^{2k-2}i}{k!\prod_{i=1}^{k-1}2i}x^k\\
           &=1+\sum_{k=1}^{\infty}-\frac{2^k(2k-2)!}{k!2^{k-1}(k-1)!}x^k\\
           &=1+\sum_{k=1}^{\infty}-\frac{2(2k-2)!}{k!(k-1)!}x^k
\ea\]
\[\ba
C(x)&=\frac{1-\sqrt{1-4x}}{2x}\\
    &=\sum_{k=1}^{\infty}\frac{(2k-2)!}{k!(k-1)!}x^{k-1}\\
    &=\sum_{n=0}^{\infty}\frac{(2n)!}{(n+1)!n!}x^n\\
    &=\sum_{n=0}^{\infty}\frac{1}{(n+1)}\binom{2n}{n}x^n
\ea\]
\[T_n=\frac{1}{(n+1)}\binom{2n}{n},\quad n\in\bbN_0\]
\epr
\ssc{AVL tree (Adelson-Velsky and Landis tree)}
\sssc{AVL tree (Adelson-Velsky and Landis tree)}
An AVL tree is a self-balancing binary search tree (BST) that keeps info of height, invented by Georgy Adelson-Velsky and Evgenii Landis. An AVL tree needs to comply with the property that for every node with two children in the tree, the height of its left and right subtrees differ by at most 1.

Some define heights of nodes and AVL trees as the definition used here plus 1, which is not used here.
\sssc{Relation of height and numbers of nodes}
Let the minimum number of nodes in an AVL tree of height $h$ be $N_h$. Since the minimum numbers of subtrees height of the root of an AVL tree of height $h$ are $h-1$ and $h-2$, we have:
\[N_{-1}=0\]
\[N_0=1\]
\[N_h=N_{h-1}+N_{h-2}+1,\quad h\in\bbN\]
\[N_h=\frac{1}{\sqrt{5}}\qty(\qty(\frac{1+\sqrt{5}}{2})^{h+3}-\qty(\frac{1-\sqrt{5}}{2})^{h+3})-1,\quad h+1\in\bbN_0\]
\bpr
Use characteristic equation:
\[x^2-x-1=0 \]
Let
\[\varphi=\frac{1+\sqrt{5}}{2}\]
Note that
\[\frac{3+\sqrt{5}}{2}=\varphi^2=\varphi+1\]
\[2+\sqrt{5}=\varphi^3=2\varphi+1\]
\[\frac{7+3\sqrt{5}}{2}=\varphi^4=3\varphi+2\]
\[\frac{1-\sqrt{5}}{2}=-\frac{1}{\varphi}=1-\varphi\]
\[\frac{3-\sqrt{5}}{2}=\frac{1}{\varphi^2}=1-\frac{1}{\varphi}=2-\varphi\]
\[x=\frac{1\pm\sqrt{5}}{2}=\varphi,-\frac{1}{\varphi}\]
\[N_h=a\varphi^h+b\qty(-\frac{1}{\varphi})^h+c\]
\[N_0=a+b+c=1\]
\[N_1=a\varphi-\frac{b}{\varphi}+c=0+1+1=2\]
\[N_2=a\varphi^2+\frac{b}{\varphi^2}+c=1+2+1=4\]
\[N_2=a\qty(\varphi+1)-b\qty(\frac{1}{\varphi}-1)+c=4\]
\[N_2-N_1=a+b=4-2=2\]
\[c=N_0-(a+b)=1-2=-1\]
\[N_1-c=a\varphi-\frac{2-a}{\varphi}=2+1=3\]
\[a\qty(\varphi^2+1)=3\varphi+2\]
\[a=\frac{3\varphi+2}{\varphi+2}=\frac{\varphi^4}{\varphi+2}\]
\[b=2-a=\frac{2-\varphi}{\varphi+2}=\frac{1}{\varphi^2\qty(\varphi+2)}\]
\[N_h=\frac{1}{\varphi+2}\qty(\varphi^{h+4}+\qty(-\frac{1}{\varphi})^{h+2})-1\]
\[\frac{1}{\varphi+2}=\frac{2}{5+\sqrt{5}}=\frac{1}{\sqrt{5}\varphi}\]
\[N_h=\frac{1}{\sqrt{5}}\qty(\qty(\frac{1+\sqrt{5}}{2})^{h+3}-\qty(\frac{1-\sqrt{5}}{2})^{h+3})-1\]
\epr
So
\[N_h=\Theta(\qty(\frac{1+\sqrt{5}}{2})^h)\]
\[h=\log_{\frac{1+\sqrt{5}}{2}}N_h+O(1)\approx 1.44\log_2N_h+O(1)\]
Let the maximum number of nodes in an AVL tree of height $h$ be $M_h$. Since the maximum numbers of subtrees height of the root of an AVL tree of height $h$ are $h-1$, we have:
\[M_{-1}=0\]
\[M_h=2M_{h-1}+1,\quad h\in\bbN_0\]
\[M_h=2^{h+1}-1,\quad h+1\in\bbN_0\]
So
\[M_h\in\Theta(\log_2M_h)\]
\[h=\log_2M_h+O(1)\]
Let the number of nodes be $n$.
\[N_h\le n\le M_h\]
\[\log_2n+O(1)\le h\le\log_{\frac{1+\sqrt{5}}{2}}n+O(1)\approx 1.44\log_2N_h+O(1)\]
Thus
\[h\in\Theta(\log n)\]
\sssc{Searching and finding minimum/maximum}
Searching and finding minimum/maximum in AVL tree is the same as in BST.

The worst-case and amortized time complexities are $O(h)=O(\log n)$.
\sssc{Rotation}
When insertion/deletion, the AVL property may be violated, called unbalanced. Since an insertion/deletion involves adding/deleting a single node, this can only increase/decrease the height of some subtree by 1. Thus, if the AVL property is violated at a node $x$, it means that the heights of left and right subtrees of it differ by exactly 2. We need to transform the tree to restore the AVL property, called rebalancing, using a single/simple rotation or a double rotation, each has two symmetric variants:
\bit
\item LL case / right single rotation: Let the left child of $x$ be $y$. The LL case is where the left subtree of $y$ has height $h+1$, while the right subtree of $y$ and the right subtree of $x$ have height $h$. We perform a right rotation on the edge between $y$ and $x$, i.e.,
\[\operatorname{left}(x)=\operatorname{right}(y)\]
\[\operatorname{right}(y)=x\]
\item RR case / left single rotation: Let the right child of $x$ be $y$. The RR case is where the right tree of $y$ has height $h+1$, while the left subtree of $y$ and the left subtree of $x$ have height $h$. We perform a left rotation on the edge between $y$ and $x$, i.e.,
\[\operatorname{right}(x)=\operatorname{left}(y)\]
\[\operatorname{left}(y)=x\]
\item LR case / left-then-right double rotation: Let the left child of $x$ be $y$ and the right child of $y$ be $z$. The LR case is where the right subtree of $y$ has height $h+1$, while the left subtree of $y$ and the right subtree of $x$ have height $h$. We perform a left rotation on the edge between $z$ and $y$ and then a right rotation on the edge between $z$, now at the original position of $y$, and $x$, i.e.,
\[\operatorname{right}(y)=\operatorname{left}(z)\]
\[\operatorname{left}(z)=y\]
\[\operatorname{left}(x)=\operatorname{right}(z)\]
\[\operatorname{right}(z)=x\]
\item RL case / right-then-left double rotation: Let the right child of $x$ be $y$ and the left child of $y$ be $z$. The RL case is where the left subtree of $y$ has height $h+1$, while the right subtree of $y$ and the left subtree of $x$ have height $h$. We perform a right rotation on the edge between $z$ and $y$ and then a left rotation on the edge between $z$, now at the original position of $y$, and $x$, i.e.,
\[\operatorname{left}(y)=\operatorname{right}(z)\]
\[\operatorname{right}(z)=y\]
\[\operatorname{right}(x)=\operatorname{left}(z)\]
\[\operatorname{left}(z)=x\]
\eit
\sssc{Insertion}
Insertion in AVL tree involves first doing insertion as with a normal BST, and then tracing the path from the new leaf to the root: For each node $x$ encountered, check if heights of its left and right subtrees differ by at most 1,
\bit
\item if yes, proceed to the parent of $x$;
\item if not, we say the AVL property is violated at $x$ and perform rotations to restore it, and then break the loop since no second violation may exist.
\eit
\bpr
Let $x$ be the first node encountered when traversing upward from the newly inserted node toward the root that violates the AVL property. Since an insertion involves adding a single node, this can only increase the height of some subtree by 1. Thus, the heights of the left and right subtrees differ by exactly 2. For any of the four cases (LL, RR, LR, RL), since the inserted node must be in the subtree of height $h+1$ of $y$, the height of the subtree whose root is $x$ has height $h+2$ both before insertion and after rotation. Thus, the fact that the AVL property is satisfied for all ancestors of $x$ before insertion implies that the AVL property is satisfied for all ancestors of $x$ after rotation.
\epr
The worst-case and amortized time complexities are $O(h)=O(\log n)$.
\sssc{Deletion}
Deletion in AVL tree involves first doing deletion the same as in BST, and then tracing the path from
\bit
\item the original parent of the deleted node if the deleted node had zero or one child, and
\item the original parent of the inorder successor or inorder predecessor that was moved if the deleted node had two children,
\eit
to the root: For each node $x$ encountered, check if heights of its left and right subtrees differ by at most 1,
\bit
\item if yes, proceed to the parent of $x$;
\item if not, we say the AVL property is violated at $x$ and perform rotations to restore it, and then proceed to the parent of $x$ since violation may propagate upwards.
\eit
The worst-case and amortized time complexities are $O(h)=O(\log n)$.
\sssc{Number of distinct AVL tree shapes}
The number $A_{h,n}$ of distinct shapes of AVL trees with $n$ nodes and height $h$ is:
\[A_{-1,0}=1\]
\[A_{-1,n}=0,\quad n\ne0\]
\[A_{0,1}=1\]
\[A_{0,n}=0,\quad n\ne1\]
\[A_{h,n}=\sum_{i=0}^{n-1}A_{h-1,i}A_{h-1,n-1-i}+A_{h-1,i}A_{h-2,n-1-i}+A_{h-2,i}A_{h-1,n-1-i},\quad h\in\bbN,n\in\bbN_0\]
\bpr
The root uses one node, leaving $n-1$ nodes for the two subtrees. Since the whole tree has height $h$, the two trees must either have heights $h-1,h-2$ or both have height $h-1$. Thus,
\[A_{h,n}=\sum_{i=0}^{n-1}A_{h-1,i}A_{h-1,n-1-i}+A_{h-1,i}A_{h-2,n-1-i}+A_{h-2,i}A_{h-1,n-1-i}\]
\epr
\sssc{Number of distinct AVL tree shapes with minimum number of nodes given height}
The number $B_h$ of distinct shapes of AVL trees with height $h$ and minimum number of nodes is:
\[B_{-1}=1\]
\[B_0=1\]
\[B_h=2B_{h-1}B_{h-2},\quad h\in\bbN\]
\[B_h=2^{\frac{1}{\sqrt{5}}\qty(\qty(\frac{1+\sqrt{5}}{2})^{h+2}-\qty(\frac{1-\sqrt{5}}{2})^{h+2})-1},\quad h+1\in\bbN_0\]
\bpr
To minimize the number of nodes, the two subtrees of root must be the AVL tree with height $h-1,h-2$ and minimum number of nodes, which have $B_{h-1},B_{h-2}$ distinct shapes respectively. There are two possible arrangements:
\bit
\item Left subtree has height $h-1$, and right subtree has height $h-2$.
\item Left subtree has height $h-2$, and right subtree has height $h-1$.
\eit
Thus,
\[B_h=2B_{h-1}B_{h-2}\]
Let
\[C_h=\log_2B_h\]
\[C_h=C_{h-1}+C_{h-2}+1\]
Use characteristic equation:
\[x^2-x-1=0\]
Let
\[\varphi=\frac{1+\sqrt{5}}{2}\]
Note that
\[\frac{3+\sqrt{5}}{2}=\varphi^2=\varphi+1\]
\[2+\sqrt{5}=\varphi^3=2\varphi+1\]
\[\frac{1-\sqrt{5}}{2}=-\frac{1}{\varphi}=1-\varphi\]
\[x=\frac{1\pm\sqrt{5}}{2}=\varphi,-\frac{1}{\varphi}\]
\[C_h=a\varphi^h+b\qty(-\frac{1}{\varphi})^h+c\]
\[C_0=a+b+c=\log_21=0\]
\[C_1=a\varphi-\frac{b}{\varphi}+c=\log_2(2*1*1)=1\]
\[C_2=a\varphi^2+\frac{b}{\varphi^2}+c=0+1+1=2\]
\[C_2=a\qty(\varphi+1)-b\qty(\frac{1}{\varphi}-1)+c=2\]
\[C_2-C_1=a+b=2-1=1\]
\[c=C_0-(a+b)=0-1=-1\]
\[C_1-c=a\varphi-\frac{1-a}{\varphi}=1+1=2\]
\[a\qty(\varphi^2+1)=2\varphi+1\]
\[a=\frac{2\varphi+1}{\varphi+2}=\frac{\varphi^3}{\varphi+2}\]
\[b=1-a=\frac{1-\varphi}{\varphi+2}=-\frac{1}{\varphi\qty(\varphi+2)}\]
\[C_h=\frac{1}{\varphi+2}\qty(\varphi^{h+3}+\qty(-\frac{1}{\varphi})^{h+1})-1\]
\[\frac{1}{\varphi+2}=\frac{2}{5+\sqrt{5}}=\frac{1}{\sqrt{5}\varphi}\]
\[C_h=\frac{1}{\sqrt{5}}\qty(\qty(\frac{1+\sqrt{5}}{2})^{h+2}-\qty(\frac{1-\sqrt{5}}{2})^{h+2})-1\]
\[B_h=2^{\frac{1}{\sqrt{5}}\qty(\qty(\frac{1+\sqrt{5}}{2})^{h+2}-\qty(\frac{1-\sqrt{5}}{2})^{h+2})-1}\]
\epr
\ssc{Splay tree}
\sssc{Splay tree}
A splay tree is a self-adjusting binary search tree (BST) that keeps no info of height and in which every access operation moves the accessed node to the root through an operation called splaying, invented by Daniel Dominic Sleator and Robert Endre Tarjan. The idea is locality, that is,  if an item is accessed, it is likely to be accessed again soon. The worst-case searching, insertion, and deletion time complexities are as with a normal BST, while the amortized time complexities are $O(\log n)$.
\sssc{(Bottom-up) splay}
When a node $x$ is accessed, a (bottom-up) splay operation consisting of a sequence of splay steps is performed on $x$ to move it to the root, called splaying $x$. Let the parent of $x$ be $p$ and the parent of $p$ be $g$. There are three types of splay steps, each of which has two symmetric variants:
\bit
\item Zig step / simple rotation: This is done if $x$ is a child of the root. We perform a rotation on the edge between $x$ and root. This only happens as the last step and only happens when $x$ originally has odd depth.
\item Zig-zig step: This is done either when $x$ is a left child of $p$ and $p$ is a left child of $g$, or when $x$ is a right child of $p$ and $p$ is a right child of $g$. We perform a rotation on the edge between $p$ and $g$ and then a rotation on the edge between $x$ and $p$.
\item Zig-zag step: This is done either when $x$ is a left child of $p$ and $p$ is a right child of $g$, or when $x$ is a right child of $p$ and $p$ is a left child of $g$. We perform a rotation on the edge between $x$ and $p$ and then a rotation on the edge between $x$, now at the original position of $p$, and $g$.
\eit
A zig step is a 1-level splay step while a zig-zig or zig-zag step is a 2-level splay step.
\sssc{Searching and finding minimum/maximum}
Searching and finding minimum/maximum in splay tree involves first doing the operation as with a normal BST and then splay the found node if found.
\sssc{Join}
Given two BSTs $S$ and $T$ such that all elements of $S$ are smaller than all elements of $T$, the following steps can be used to join them to a single BST: either
\ben
\item Splay the largest item in $S$. Now this item is in the root $r$ of $S$ and has no right child.
\item Set the right child of the new root $r$ of $S$ to the root of $T$.
\een
or
\ben
\item Splay the smallest item in $T$. Now this item is in the root $r$ of $T$ and has no left child.
\item Set the left child of the new root $r$ of $T$ to the root of $S$.
\een
\sssc{Split}
Given a BST and a value $x$, splitting the BST at $x$ returns two new BSTs, one containing all elements less than or equal to $x$, $L$, and the other containing all elements greater than $x$, $R$. This can be done with:
\ben
\item Search for $x$ in the tree. If we found it, let $y$ be that node; otherwise let $y$ be the node we stop at.
\item Splay $y$.
\item Split its right subtree, which is $R$, from the rest of the tree, which is $L$.
\een
\sssc{Insertion}
To insert a value $x$ into a splay tree, there are two ways: either
\ben
\item Insert $x$ as with a normal BST.
\item Splay $x$.
\een
or
\ben
\item Split at $x$. However, if $x$ is found in the tree, stop after splaying $x$.
\item Create a new tree with $x$ as root and the two split trees as its subtrees.
\een
\sssc{Deletion}
To delete a node $x$ from a splay tree: either
\ben
\item Search for $x$.
\item Splay $x$.
\item Delete $x$ and join the two subtrees of $x$.
\een
or
\ben
\item Delete $x$ as with a normal BST.
\item Splay the original parent of the deleted $x$.
\een
\sssc{Amortized analysis}
We define size of a node $s(x)$ as the number of nodes in the subtree whose root is node $x$, rank of a node as
\[r(x)=\log_2s(x)\]
and potential function of a tree $P(T)$ is
\[P(T)=\sum_{x\in T}r(x)\]
The actual cost $C_i$ of an operation $i$ is defined as the number of rotations in the operation. Let the state before operation $i$ be denoted by the original functions and the state after operation $i$ be denoted with a prime symbol. The amortized cost $A_i$ of an operation $i$ with actual cost $C_i$ that transforms the tree $T_{i-1}$ to $T_i$ is defined as
\[A_i=C_i+\Delta P=C_i+P(T_i)-P(T_{i-1})\]
\bit
\item Zig step:
\[C_i=1\]
\[r'(p)\le r(p)\]
\[\Delta P=r'(x)+r'(p)-r(x)-r(p)\le r'(x)-r(x)\]
\[A_i\le 1+r'(x)-r(x)\]
\item Zig-zig step:
\[C_i=2\]
\[r'(p)\le r(g)\]
\[r(x)\le r(p)\]
\[s(x)+s'(g)=s'(x)-2\le s'(x)\]
The log function is concave, i.e.,
\[\frac{\log_2a+\log_2b}{2}\le\log_2(\frac{a+b}{2})\]
i.e.,
\[log_2a+\log_2b\le 2\log_2(a+b)-2\]
Hence
\[r(x)+r'(g)=2r'(x)-2\]
\[r'(g)\le 2r'(x)-r(x)-2\]
\[\Delta P=r'(x)+r'(p)+r'(g)-r(x)-r(p)-r(g)\le r'(x)+r'(g)-2r(x)\le 3(r'(x)-r(x))-2\]
\[A_i\le 3(r'(x)-r(x))\]
\item Zig-zag step:
\[C_i=2\]
\[s'(p)+s'(g)=s'(x)-1\le s'(x)\]
The log function is concave, i.e.,
\[\frac{\log_2a+\log_2b}{2}\le\log_2(\frac{a+b}{2})\]
i.e.,
\[log_2a+\log_2b\le 2\log_2(a+b)-2\]
Hence
\[r'(p)+r'(g)\le 2r'(x)-2\]
\[r(x)\le r(p)\le r(g)\]
\[\Delta P=r'(x)+r'(p)+r'(g)-r(x)-r(p)-r(g)\le 3(r'(x)-r(x))-2\]
\[A_i\le 3(r'(x)-r(x))\]
\eit
When splaying a node $x$, zero or more 2-level splay steps followed by zero or one 1-level splay step is conducted. Let $r_2(x)$ be the rank of $x$ just after the last 2-level splay step and $r_f(x)$ be the final rank of $x$. The amortized cost $A$ of splaying $x$ is
\[\ba
A&\le 1+r_f(x)-r_2(x)+3(r_2(x)-r(x))\\
 &\le 1+3(r_f(x)-r_2(x))+3(r_2(x)-r(x))\\
 &=1+3(r_f(x)-r(x))\\
 &\le 1+3(\log_2n-r(x))\\
 &=O(\log n)
\ea\]
For a sequence of $m$ accesses in a splay tree $T$, we can bound the total actual cost with
\[\sum_{i=1}^mC_i=P(T_0)-P(T_m)+\sum_{i=1}^mA_i\]
For every node $x$,
\[1\le s(x)\le n\]
\[0\le r(x)\le\log_2n\]
Thus the change of rank $\Delta r(x)$ of a node $x$ is
\[\Delta r(x)\le\log_2n\]
Thus the change of potential function is
\[P(T_0)-P(T_m)\le\sum_{x\in T}\log_2n=O\qty(n\log n)\]
\sssc{Weighted analysis}
The above analysis can be generalized by assigning each node $x$ a weight $w(x)\ge 1$ and defining $s(x)$ as the sum of weights of nodes in the subtree whose root is node $x$, called weighted analysis. Note that the previous analysis corresponds to the special case $w(x)=1$. The sum of all weights is denoted $W$. The same analysis gives:
\[A\le 1+3(\log_2n-r(x))=O(\log(\frac{W}{w(x)})+1)\]
For a sequence of $m$ accesses in a splay tree $T$, we can bound the total actual cost with
\[\sum_{i=1}^mC_i=P(T_0)-P(T_m)+\sum_{i=1}^mA_i\]
For every node $x$,
\[w(x)\le s(x)\le W\]
\[\log_2(w(x))\le r(x)\le\log_2W\]
Thus the change of rank $\Delta r(x)$ of a node $x$ is
\[\Delta r(x)\le\log_2(\frac{W}{w(x)})\]
Thus the change of potential function is
\[P(T_0)-P(T_m)\le\sum_{x\in T}\log_2(\frac{W}{w(x)})=O\qty(\sum_{x\in T}\log(\frac{W}{w(x)}))\]
\sssc{Balance theorem}
The cost of performing a sequence of $m$ accesses in a splay tree containing $n$ nodes is
\[O(m\log n+n\log n)\]
\bpr
Take a constant weight, e.g.,
\[w(x)=1\]
Then
\[W=n\]
\[A_i=O(\log n)\]
\[P(T_0)-P(T_m)=O(n\log n)\]
\epr
This theorem implies that splay trees perform at least as well as AVL trees on sequences of at least $n$ accesses.
\sssc{Static optimality theorem}
The cost of performing a sequence of $m$ accesses in a splay tree $T$ containing $n$ nodes, with each node accessed at least once and $q_x$ being the number of times node $x$ is accessed, is
\[O\qty(m+\sum_{x\in T}q_x\log(\frac{m}{q_x}))\]
\bpr
Let
\[w(x)=q_x\]
Then
\[W=m\]
\[A_i=O\qty(\log(\frac{m}{q_x})+1)\]
\[\sum_{i=1}^mA_i=O\qty(\sum_{x\in T}q_x\log(\frac{m}{q_x})+m)\]
\[q_x\ge 1\]
\[P(T_0)-P(T_m)=O\qty(\sum_{x\in T}\log(\frac{m}{q_x}))=O\qty(\sum_{x\in T}q_x\log(\frac{m}{q_x}))\]
\epr
\sssc{Working set theorem}
The cost of performing a sequence $S$ of $m$ accesses in a splay tree $T$ containing $n$ nodes, with $t(x)$ being the number of distinct nodes accessed before the previous time node $x$ was accessed for each $x\in S$, is
\[O\qty(m+n\log n+\sum_{x\in S}\log(t(x)+1))\]
\bpr
Let
\[w(x)=(t(x)+1)^{-2}\]
Then
\[W=O(1)\]
\[A_i=O\qty(\log(t(x)+1)+1)\]
\[\sum_{i=1}^mA_i=O\qty(\sum_{x\in S}\log(t(x)+1)+m)\]
\[P(T_0)-P(T_m)=O\qty(n\log n)\]
\epr
\sssc{Locality property or strong balance theorem}
The cost of performing a sequence of $m$ accesses in a splay tree containing $n$ nodes, with $k$ distinct nodes accessed, is
\[O(m+m\log k+n\log n)\]
\bpr
Use working set theorem.
\[t(x)=O(k)\]
\[O\qty(m+n\log n+\sum_{x\in S}\log(t(x)+1))=O(m+n\log n+m\log k)\]
\epr
\sssc{Top-down splay}
An alternative implementation of splay that is done together with searching. Suppose we want to splay $x$. We start at the root and maintain two trees which are initially empty, left tree $L$ of nodes smaller than $x$ and right tree $R$ of nodes larger than $x$, during the traversal. In each step with current node $c$, there are three cases, each with two symmetric variants, where we will describe the former and the latter is the former with left and right exchanged and $L$ and $R$ exchanged:
\bit
\item $x<\operatorname{left}(c)$: We perform a rotation on the edge between $c$ and $\operatorname{left}(c)$ and then a rotation one the edge $\operatorname{left}(\operatorname{left}(c))$ and $\operatorname{left}(c)$, and then assign $\operatorname{left}(c)$ to $c$.
\item $x>\operatorname{right}(c)$: We perform a rotation on the edge between $c$ and $\operatorname{right}(c)$ and then a rotation one the edge $\operatorname{right}(\operatorname{right}(c))$ and $\operatorname{right}(c)$, and then assign $\operatorname{right}(c)$ to $c$.
\item $\operatorname{left}(c)<x<c$: We connect $c$ into $R$, and then assign $\operatorname{left}(c)$ to $c$.
\item $\operatorname{right}(c)>x>c$: We connect $c$ into $L$, and then assign $\operatorname{right}(c)$ to $c$.
\item $x=c$, $x<c$ and $c$ has no left child, or $x>c$ and $c$ has no right child: The traversal stops at $c$. We take $c$ as the new root, and $L$ and $R$ as left and right subtrees of $c$.
\eit
Top-down splay has the same worst-case and amortized searching, insertion, and deletion time complexities as bottom-up splay and $O(1)$ extra space.
\ssc{An implementation of BST, ASL tree, and splay tree in Python}
\lstinputlisting[language=Python]{bst.py}
\end{document}
